\section{Evaluation}
\label{sec:eval}

\subsection{Setup}
\label{sec:eval-setup}

Each engine is synthesised out of context with Yosys~\cite{wolf_yosys} through
the OpenROAD flow~\cite{ajayi2019openroad} at the Nangate45
\SI{45}{\nano\meter} node, under a common timing constraint.
Functional simulation uses Verilator~\cite{snyder_verilator}.

\subsection{Synthesis and throughput analysis}
\label{sec:eval-synthesis}

Physical synthesis results confirm our core hypothesis: specialising hardware datapaths to narrow microscaling formats yields significant area, frequency, and throughput density advantages over wider uniform representations.

First, our measurements demonstrate that narrow fused datapaths achieve superior physical efficiency. \fpfour{} emerges as the most compact and agile fused engine, requiring only \SI{13442}{\micro\meter\squared} and operating at \SI{414.9}{\mega\hertz}. This efficiency stems directly from its narrow four-bit multipliers and lightweight accumulation frame, compared to the 95-bit frame required by \fpeight. While \mtwo{} offers metadata refinements, its extra logic increases area and power consumption to \SI{441.0}{\milli\watt} while lowering maximum frequency to \SI{359.5}{\mega\hertz}. \fpeight{} requires the largest area (\SI{23881}{\micro\meter\squared}) and lowest frequency (\SI{288.3}{\mega\hertz}), confirming that wide uniform formats incur severe area and timing penalties that offset their element-level parallel gains.

Second, we establish that multi-instruction residue execution can overcome register file port constraints without sacrificing compute throughput. By decomposing the \fpfourres{} dot product into a decoupled instruction pair (\insn{mxr4dot} and \insn{mxfinal}), we eliminate the three-read-port register file bottleneck while unrolling combinational delay paths. The reduction stage (\insn{mxr4dot}) achieves \SI{915.5}{\mega\hertz}, and the accumulator stage (\insn{mxfinal}) reaches \SI{441.0}{\mega\hertz}. Sustaining one completed pair every two cycles yields a useful throughput of \SI{14.11}{\giga\flop\per\second}, representing a \SI{3.06}{\times} speedup over \fpeight{} (Fig.~\ref{fig:throughput_comparison}) and outperforming all single-instruction candidates.



\subsection{Unified evaluation framework}
\label{sec:eval-compare}

Because prior studies evaluated microscaling formats across different process nodes, synthesis flows, and host processors, published figures cannot be directly compared to isolate format-level trade-offs. To establish a rigorous benchmark, we reimplemented all format engines within a unified framework using the same RISC-V host core (\core), identical timing constraints, the open-source Yosys and OpenROAD toolchain, and the Nangate45 \SI{45}{\nano\meter} PDK.